% FaceCheck 设计文档
% 编译: xelatex design.tex

\documentclass[11pt,a4paper]{ctexart}

% ========== 页面设置 ==========
\usepackage[top=2.5cm,bottom=2.5cm,left=2.8cm,right=2.8cm]{geometry}

% ========== 引擎检测 ==========
\usepackage{iftex}
\ifPDFTeX
  % pdflatex: 需要打补丁让 TikZ 节点正确渲染中文
  \usepackage{etoolbox}
  \makeatletter
  \patchcmd{\tikz@installcommands}{\tikz@node@textfont}
    {\tikz@node@textfont\CJKfamily{zhrm}}{}{}
  \makeatother
\fi

% ========== 图表与图形 ==========
\usepackage{tikz}
\usetikzlibrary{shapes.geometric, arrows.meta, positioning, fit, backgrounds, calc}

% ========== 列表增强 ==========
\usepackage{enumitem}
\setlist{nosep, leftmargin=2em}

% ========== 超链接 ==========
\usepackage[colorlinks=true,linkcolor=black,urlcolor=blue]{hyperref}

% ========== 代码等宽字体 ==========
\usepackage{inconsolata}
\usepackage{xcolor}

% ========== 表格 ==========
\usepackage{array}
\usepackage{booktabs}
\usepackage{tabularx}

% ========== 页眉页脚 ==========
\usepackage{fancyhdr}
\pagestyle{fancy}
\fancyhf{}
\fancyhead[L]{FaceCheck 设计文档}
\fancyhead[R]{\thepage}
\renewcommand{\headrulewidth}{0.4pt}

% ========== 自定义颜色 ==========
\definecolor{springboot}{HTML}{6DB33F}
\definecolor{pgcolor}{HTML}{336791}
\definecolor{rediscolor}{HTML}{DC382D}
\definecolor{rabbitcolor}{HTML}{FF6600}
\definecolor{frscolor}{HTML}{0078D4}
\definecolor{obscolor}{HTML}{FFB900}
\definecolor{fluttercolor}{HTML}{02569B}
\definecolor{boxbg}{HTML}{F5F5F5}
\definecolor{border}{HTML}{666666}

% ========== 自定义命令 ==========
\newcommand{\tech}[1]{\texttt{\small#1}}

\begin{document}

\title{\Huge FaceCheck 设计文档}
\author{管理员场次人脸签到系统}
\date{\today}
\maketitle
\thispagestyle{fancy}

\tableofcontents
\newpage

% ======================================================================
\section{目标与范围}

\begin{itemize}[leftmargin=2.5em]
  \item 交付完整的人脸签到系统：管理员场次管理、匿名扫码签到、个人照片库、个人与全局签到记录。
  \item 后端保持 Spring Boot 单体作为唯一业务边界。
  \item 人脸识别使用华为云 FRS，图片存储使用华为云 OBS。
  \item 匿名用户、普通用户、管理员三类访问边界清晰隔离。
\end{itemize}

% ======================================================================
\section{非目标}

\begin{itemize}[leftmargin=2.5em]
  \item 不拆分微服务。
  \item 不引入本地 OpenCV/ONNX 推理。
  \item 客户端不直连 FRS/OBS/PostgreSQL/Redis/RabbitMQ。
  \item 第一阶段不做活体检测。
\end{itemize}

% ======================================================================
\section{架构概览}

\subsection{系统上下文}

\begin{figure}[htbp]
\centering
\begin{tikzpicture}[
    node distance=0.6cm and 0.5cm,
    box/.style={
        rectangle, draw=border, fill=boxbg, rounded corners=3pt,
        minimum width=3.8cm, minimum height=0.65cm, align=center, font=\footnotesize
    },
    appbox/.style={
        box, fill=fluttercolor!10, draw=fluttercolor!50!black, minimum width=3cm
    },
    modbox/.style={
        rectangle, draw=springboot!30, fill=springboot!3,
        rounded corners=3pt, inner sep=6pt, font=\footnotesize, align=left,
        minimum width=4.2cm, minimum height=2.2cm
    },
    extbox/.style={
        box, minimum width=4.2cm, minimum height=0.7cm, font=\footnotesize
    },
    arrow/.style={->, >=Stealth, thick, draw=border},
    darrow/.style={->, >=Stealth, thick, dashed, draw=border},
]

% ====== Flutter ======
\node[appbox] (flutter) {Flutter App};

% ====== Module nodes (placed inside backend scope) ======
\node[modbox, below=1.2cm of flutter] (bizmod) {
    \textbf{业务模块}\\[2pt]
    auth \quad identity \quad session\\
    checkin \quad admin
};

\node[modbox, right=0.6cm of bizmod] (infmod) {
    \textbf{基础设施模块}\\[2pt]
    face \quad (FaceRecognitionProvider)\\
    storage \quad (OBS 适配)\\
    infrastructure \quad (安全/日志/SDK/消息)
};

% ====== Fit backend box around modules ======
\node[rectangle, draw=springboot!50!black, fill=none,
      rounded corners=5pt, inner sep=8pt,
      fit=(bizmod)(infmod), label=above:\footnotesize\textbf{Spring Boot 单体}] (backend) {};

% ====== Middle row: DB / Cache / MQ ======
\node[extbox, below=1.0cm of backend,
      fill=pgcolor!10, draw=pgcolor!50!black] (pg)
      {PostgreSQL (事实源)};

\node[extbox, left=0.4cm of pg,
      fill=rediscolor!10, draw=rediscolor!50!black] (redis)
      {Redis (限流/幂等/黑名单)};

\node[extbox, right=0.4cm of pg,
      fill=rabbitcolor!10, draw=rabbitcolor!50!black] (rabbit)
      {RabbitMQ (注册/删除补偿)};

% ====== Bottom row: Huawei Cloud ======
\node[extbox, below=0.6cm of redis,
      fill=frscolor!10, draw=frscolor!50!black] (frs)
      {Huawei Cloud FRS};

\node[extbox, below=0.6cm of rabbit,
      fill=obscolor!10, draw=obscolor!50!black] (obs)
      {Huawei Cloud OBS};

% ====== Arrows ======
\draw[arrow] (flutter) -- (backend);

\draw[arrow] (backend.south -| redis.north) -- (redis.north);
\draw[arrow] (backend.south -| pg.north)    -- (pg.north);
\draw[arrow] (backend.south -| rabbit.north) -- (rabbit.north);

\draw[darrow] (redis.south)  -- (frs.north);
\draw[darrow] (rabbit.south) -- (obs.north);

\end{tikzpicture}
\caption{系统上下文架构图}
\label{fig:architecture}
\end{figure}

\subsection{关键边界}
\begin{itemize}[leftmargin=2.5em]
  \item App 只调用后端 API。
  \item FRS/OBS 只由后端访问。
  \item PostgreSQL 是唯一业务事实源。
\end{itemize}

% ======================================================================
\section{模型选择}

\subsection{人脸识别模型}
\begin{itemize}[leftmargin=2.5em]
  \item \textbf{提供方：}华为云 FRS。
  \item \textbf{调用链路：} \texttt{detect} $\rightarrow$ \texttt{search} $\rightarrow$ \texttt{compare}。
  \item \textbf{抽象接口：} \tech{FaceRecognitionProvider}，生产实现为 \tech{HuaweiCloud\-Frs\-Face\-Recognition\-Provider}，测试实现为 \tech{MockFaceRecognitionProvider}。
  \item \textbf{理由：}符合约束，避免本地推理，便于替换与测试隔离。
\end{itemize}

\subsection{图片存储}
\begin{itemize}[leftmargin=2.5em]
  \item \textbf{提供方：}华为云 OBS。
  \item \textbf{对象路径：}
  \begin{itemize}
    \item 人脸照片：\texttt{faces/user/\{userId\}/\{photoId\}.jpg}
    \item 签到照片：\texttt{checkins/session/\{sessionId\}/attempt/\{attemptId\}.jpg}
  \end{itemize}
  \item \textbf{理由：}对象存储适合二进制图片，元数据与业务关系仍由 PostgreSQL 维护。
\end{itemize}

\subsection{数据与消息}
\begin{itemize}[leftmargin=2.5em]
  \item \textbf{PostgreSQL：}用户、场次、照片、attempt、record、审计日志、外部调用日志的权威存储。
  \item \textbf{Redis：}Token 黑名单、限流、幂等结果缓存。
  \item \textbf{RabbitMQ：}人脸照片异步注册与删除补偿，各含重试与死信队列。
\end{itemize}

% ======================================================================
\section{关键数据模型}

第一阶段核心实体：

\begin{table}[htbp]
\centering
\caption{核心实体概览}
\label{tab:entities}
\footnotesize
\begin{tabularx}{\textwidth}{@{}>{\raggedright\arraybackslash}p{2.6cm} >{\raggedright\arraybackslash}X >{\raggedright\arraybackslash}X@{}}
\toprule
实体 & 关键字段 & 约束/状态 \\
\midrule
User & username, password\_hash, role, status, lastLoginAt &
  表名 \texttt{user\_account}\newline ROLE: \texttt{USER|ADMIN}\newline STATUS: \texttt{ACTIVE|LOCKED|DISABLED} \\

FacePhoto & obsBucket, obsRegion, obsObjectKey, contentType,
           sizeBytes, sha256, storageProvider, detectStatus,
           registerStatus, failureReason, failureCode &
  DETECT: \texttt{PENDING|PASSED|FAILED}\newline REGISTER: \texttt{PENDING|ACTIVE|FAILED}\newline \texttt{DELETE\_PENDING}\newline \texttt{DELETE\_FAILED|DELETED} \\

HuaweiFaceRef & userId, facePhotoId, faceSetName, frsFaceId,
                externalImageId, externalFields(JSONB), status &
  STATUS: \texttt{ACTIVE|DELETE\_PENDING}\newline \texttt{DELETE\_FAILED}\newline \texttt{DELETED|ORPHANED} \\

Attendance\newline Session & name, description, startTime, endTime,
                     lateAfterTime, status, qrToken &
  STATUS: \texttt{DRAFT|PUBLISHED}\newline \texttt{CLOSED|CANCELED}\newline qrToken 高熵唯一 \\

Attendance\newline Checkin\newline Attempt & obsBucket, obsRegion, obsObjectKey,
                            contentType, sizeBytes, sha256,
                            status, resultCode, failureReason,
                            frsRequestId, matchedUserId, similarity,
                            idempotencyKey, reviewed &
  STATUS: \texttt{PROCESSING|SUCCESS|FAILED}\newline \texttt{DUPLICATE\_CHECKIN}，\newline
  UK(session\_id, idempotency\_key) \\

Attendance\newline Record & sessionId, userId, attemptId, checkinTime,
                   status, similarity, source &
  STATUS: \texttt{VALID|MANUAL\_CONFIRMED}\newline SRC: \texttt{APP\_QR\_ANON|ADMIN\_RETRY}，\newline
  UK(session\_id, user\_id) + UK(attempt\_id) \\

AuditLog & actorUserId, actorRole, action, targetType,
          targetId, summary, detailJson(JSONB) & 管理员关键操作审计 \\

ExternalService\newline CallLog & serviceName, operation, requestId,
                          status, latencyMs, errorCode &
  SVC: \texttt{HUAWEI\_FRS|HUAWEI\_OBS}\newline STATUS: \texttt{SUCCESS|FAILED|TIMEOUT}\newline \texttt{RATE\_LIMITED|RETRYING} \\

SystemConfig & configKey, configValue, valueType, description &
  TYPE: \texttt{STRING|NUMBER|BOOLEAN|JSON}\newline 白名单配置管理 \\

\bottomrule
\end{tabularx}
\end{table}

\subsection{关键约束}
\begin{itemize}[leftmargin=2.5em]
  \item \tech{AttendanceRecord} 唯一约束 \texttt{unique(session\_id, user\_id)} 和 \texttt{unique(attempt\_id)}。
  \item \tech{AttendanceCheckinAttempt} 唯一约束 \texttt{unique(session\_id, idempotency\_key)}。
  \item \tech{qrToken} 高熵且唯一。
\end{itemize}

% ======================================================================
\section{关键流程}

\subsection{人脸照片上传与异步注册}
\begin{enumerate}[leftmargin=2.5em]
  \item 登录用户或管理员上传照片。
  \item 后端校验大小、MIME、扩展名、可解码性（$\le$10MB，JPEG/PNG/WEBP）。
  \item 后端强制每人最多 5 张可用照片。
  \item 上传到 OBS，创建 FacePhoto（detectStatus/registerStatus 均为 PENDING）。
  \item 发送 \tech{FacePhotoRegisterTask} 到 RabbitMQ。
  \item 消费者通过 \tech{FaceRecognitionProvider} 完成 detect + enroll。
  \item 成功写 \tech{HuaweiFaceRef} 并标记 FacePhoto ACTIVE；失败写入原因。
\end{enumerate}

\subsection{匿名扫码签到}
\begin{enumerate}[leftmargin=2.5em]
  \item 管理员发布场次并生成 qrToken。
  \item App 扫码调用公开入口确认场次可用。
  \item 提交签到照片与 idempotencyKey、deviceId。
  \item 后端限流与幂等检查。
  \item 生成 \tech{AttendanceCheckinAttempt} 并上传签到照片。
  \item 执行 \texttt{detect} $\rightarrow$ \texttt{search} $\rightarrow$ \texttt{compare}。
  \item 映射候选到 userId 并按阈值判断。
  \item 成功写 \tech{AttendanceRecord}；唯一约束冲突返回 DUPLICATE\_CHECKIN。
  \item 返回 SUCCESS / FAILED / PROCESSING / DUPLICATE\_CHECKIN。
\end{enumerate}

\subsection{场次生命周期}
\begin{itemize}[leftmargin=2.5em]
  \item \texttt{DRAFT} $\rightarrow$ \texttt{PUBLISHED} $\rightarrow$ \texttt{CLOSED}，或 \texttt{DRAFT/PUBLISHED} $\rightarrow$ \texttt{CANCELED}。
  \item qrToken 轮换后旧 token 失效。
\end{itemize}

% ======================================================================
\section{安全与权限}
\begin{itemize}[leftmargin=2.5em]
  \item JWT 保护需要登录的接口，注销写入 Redis 黑名单。
  \item 普通用户只能访问自身数据。
  \item 管理员接口仅限管理员角色。
  \item 匿名接口仅限场次解析、签到提交与结果查询。
\end{itemize}

% ======================================================================
\section{结果与错误码}

\subsection{主状态}
\begin{center}
\texttt{SUCCESS} \quad | \quad
\texttt{FAILED} \quad | \quad
\texttt{PROCESSING} \quad | \quad
\texttt{DUPLICATE\_CHECKIN}
\end{center}

\subsection{关键结果码}

\begin{table}[htbp]
\centering
\caption{签到结果码}
\label{tab:result-codes}
\begin{tabularx}{\textwidth}{@{}p{2.6cm} >{\raggedright\arraybackslash}X@{}}
\toprule
类别 & 结果码 \\
\midrule
图片问题 & \texttt{INVALID\_IMAGE} \quad \texttt{NO\_FACE} \quad \texttt{MULTIPLE\_FACES} \\
匹配结果 & \texttt{LOW\_CONFIDENCE} \quad \texttt{NO\_MATCH} \quad \texttt{DUPLICATE\_CHECKIN} \quad \texttt{MANUAL\_REVIEW} \\
服务异常 & \texttt{FRS\_TIMEOUT} \quad \texttt{FRS\_RATE\_LIMITED} \quad \texttt{FRS\_ERROR} \\
\bottomrule
\end{tabularx}
\end{table}

% ======================================================================
\section{可观测性与审计}
\begin{itemize}[leftmargin=2.5em]
  \item 请求链路统一注入 traceId。
  \item \tech{AuditLog} 记录管理员关键操作。
  \item \tech{ExternalServiceCallLog} 记录 FRS/OBS 调用 requestId 与耗时。
\end{itemize}

% ======================================================================
\section{测试策略}
\begin{itemize}[leftmargin=2.5em]
  \item \tech{MockFaceRecognitionProvider} 覆盖单元与集成测试。
  \item Testcontainers 覆盖 PostgreSQL、Redis、RabbitMQ。
  \item 关键 API 合同测试与集成测试。
  \item Android 模拟器或真机验证扫码、拍照、上传与匿名签到流程。
\end{itemize}

% ======================================================================
\section{风险与缓解}
\begin{itemize}[leftmargin=2.5em]
  \item \textbf{FRS 不稳定：}用明确结果码和 attempt 记录可追溯失败。
  \item \textbf{重复提交：}幂等缓存 + 数据库唯一约束双重保护。
  \item \textbf{照片质量：}上传阶段前置校验并输出明确失败原因。
\end{itemize}

% ======================================================================
\section{后续扩展}
\begin{itemize}[leftmargin=2.5em]
  \item 可切换为异步签到处理，但保持 PROCESSING 与结果查询契约不变。
  \item 系统状态/配置页面作为第二阶段扩展。
\end{itemize}

\end{document}
